\documentclass[conference]{IEEEtran}
\usepackage{amsmath,amssymb}
\usepackage{booktabs}
\usepackage{url}

\title{Counterfactual Red-Teaming of Autonomous Threat-Response Agents in Industrial Cyber-Physical Systems}

\author{\IEEEauthorblockN{Ichrak Hamdi}
\IEEEauthorblockA{Department of Computer Science and Software Engineering\\
Universit\'e Laval, Qu\'ebec, Canada\\
ichrak.hamdi.1@ulaval.ca}}

\begin{document}
\maketitle

\begin{abstract}
An explanation of an autonomous industrial response is operationally useful
only if the alternative it describes can occur in the defended system.
Feature-space counterfactual explanations may change detector scores, process
measurements, or network events independently, even though those quantities are
linked by attack prerequisites and process dynamics. This paper formulates
counterfactual explanation as campaign-level red-teaming of a frozen threat-response
policy. The proposed black-box search modifies the timing, visibility, duration,
and intensity of fixed multi-stage campaigns while an independent validator
enforces network reachability, privilege, temporal, and physical constraints.
Every candidate is re-executed in a closed-loop three-stage industrial process.
The returned explanation is the lowest-cost executable campaign that causes a
weaker response and a larger physical impact. The evaluation compares random,
feature-level, cyber-valid, and cyber-physical counterfactual generation and
tests whether explanations remain valid under a higher-fidelity execution
configuration. The accompanying engineering pilot verifies the complete
execution path; the confirmatory multi-seed evaluation protocol is fixed before
the reported comparison is produced.
\end{abstract}

\begin{IEEEkeywords}
explainable artificial intelligence, counterfactual explanation, autonomous
cyber defence, industrial cyber-physical systems, reinforcement learning
\end{IEEEkeywords}

\section{Introduction}

Autonomous cyber-defence agents shorten the interval between an alert and a
response, a capability that matters when an attack can propagate through a
controller into a physical process~\cite{kott2019aica}. Their decisions are also
difficult to inspect. An operator may need to know why an agent isolated an
engineering workstation, activated a backup controller, or continued to
monitor. Explainable reinforcement learning addresses this need through feature,
policy, trajectory, and causal explanations~\cite{milani2024xrl,madumal2020causal}.

Industrial response places an additional requirement on a counterfactual. A
statement such as ``the response would change if the historian alarm were 0.2
lower'' need not correspond to any possible campaign. The alarm may be fixed by
the traffic generated during historian access; the access may require an
engineering credential; and changing a measurement while holding the related
flow and tank levels fixed may violate the process dynamics. Such an explanation
is valid for the policy function but invalid for the system in which the policy
acts.

This paper uses counterfactual XAI as an offline red-teaming mechanism. Starting
from a fixed multi-stage campaign and a frozen response policy, it asks which
smallest executable change to the campaign causes a weaker response and a worse
physical outcome. The response policy is queried as a black box, while candidate
campaigns are checked independently and re-executed in closed loop.

The contributions are:
\begin{enumerate}
    \item a campaign-level formulation of counterfactual explanation for
    sequential autonomous industrial response;
    \item a black-box search whose edits obey attack-chain, network, temporal,
    and process constraints; and
    \item a closed-loop protocol measuring decision validity, campaign
    feasibility, physical impact, edit cost, and transfer validity.
\end{enumerate}

\section{Related Work and Gap}

Counterfactual explanations identify a small change under which a model would
return a different decision~\cite{wachter2017counterfactual}. In sequential
decision-making, the object being explained may be an action or trajectory, and
the relevant evidence can extend over time~\cite{milani2024xrl}. Causal models
can distinguish intervention from observation and provide a stronger basis for
explaining an RL decision~\cite{madumal2020causal}.

Recent autonomous-defence systems increasingly combine temporal reasoning and
explanation. DeepXplain provides structural, temporal, and policy-level
explanations for stage-aware APT defence~\cite{phan2026deepxplain}. Zhang et
al.~\cite{zhang2026explainable} combine a structural causal model, constrained
response transitions, and adversarial multi-agent policies. These methods
advance policy transparency, but do not establish that a counterfactual change
returned to the analyst corresponds to an executable industrial campaign with
a reproduced physical consequence.

The distinction is material in industrial systems. Stealthy attacks are already
designed under detector and process constraints~\cite{urbina2016stealthy}, and
concealment attacks explicitly manipulate evidence while preserving the attack
objective~\cite{erba2020concealment}. A useful explanation must therefore alter
the causes of the evidence through permitted tactics, not edit each observed
feature independently. Faithful environment interfaces and the separation of
hidden environment state from agent observations are also central to credible
autonomous-defence evaluation~\cite{hicks2026environments}.

\section{Problem Formulation}

Let a fixed campaign be a temporally ordered sequence
$\mathcal{C}=(e_1,\ldots,e_m)$. Event $e_j$ contains an attack stage, execution
time, asset, magnitude, visibility, and duration. Executing $\mathcal{C}$ in an
industrial environment produces an observation history $H_t$, a response
$a_t=\pi(H_t)$ from a frozen policy $\pi$, and a physical trajectory $x_{0:T}$.

The edit vector $\Delta$ changes permitted campaign fields and produces
$\mathcal{C}^{\Delta}$. We seek
\begin{equation}
\begin{aligned}
\Delta^\star = \arg\min_{\Delta}\quad & C(\Delta)\\
\text{s.t.}\quad
& \pi(H_{t_a}^{\Delta}) < \pi(H_{t_a}),\\
& J_{\mathrm{phys}}(\mathcal{C}^{\Delta})
  -J_{\mathrm{phys}}(\mathcal{C}) \geq \eta,\\
& \mathcal{C}^{\Delta}\in
  \mathcal{F}_{\mathrm{atk}}\cap\mathcal{F}_{\mathrm{net}}
  \cap\mathcal{F}_{\mathrm{time}}\cap\mathcal{F}_{\mathrm{phys}}.
\end{aligned}
\label{eq:cf}
\end{equation}
Responses are ordered by intervention strength, $t_a$ is the declared audit
instant before physical impact, and $J_{\mathrm{phys}}$ is maximum departure
from the nominal operating point after impact. The feasibility sets represent
attack prerequisites, network reachability and privilege, temporal order, and
closed-loop physical execution. The edit cost is
\begin{equation}
C(\Delta)=\sum_j w_v|\Delta v_j|+w_t|\Delta t_j|/T
 +w_d|\Delta d_j|/T+w_m|\Delta m_j|,
\end{equation}
where $v,t,d,m$ denote visibility, time, duration, and magnitude.

Equation~\eqref{eq:cf} makes the explanation adversarial in purpose but
diagnostic in use: it identifies a feasible weakness without changing or
retraining the frozen responder.

\section{Campaign-Constrained Counterfactual Search}

\subsection{Campaign Grammar}

The prototype campaign progresses through credential access, engineering access,
PLC discovery, historian impairment, measurement replay, and actuator override.
Events are aligned with ATT\&CK for ICS~\cite{mitreics}. Each stage declares its
permitted assets and prerequisite stages. For example, PLC discovery requires
prior engineering access, and measurement replay requires prior historian
impairment. The physical impact event is held fixed during explanation; the
search modifies only the evidence-producing preparation path.

Four checks are applied. The attack check requires every prerequisite. The
network check requires the event's asset to be reachable with the acquired
privilege. The temporal check enforces prerequisite order and finite durations.
The physical check executes the candidate and rejects undefined or non-finite
trajectories. The first three operate before simulation; the fourth cannot be
replaced by a feature-range check because physically coupled variables evolve
together.

\subsection{Black-Box Search}

The method uses a bounded beam search. At every depth it generates local changes
to one campaign event, validates the resulting campaign, executes it with the
frozen responder, and ranks it by response reduction, impact increase, and edit
cost. Duplicate campaign signatures are evaluated once. Search terminates at a
fixed depth, and the least-cost candidate satisfying~\eqref{eq:cf} is revalidated
and replayed before reporting.

The policy interface contains only \texttt{reset} and \texttt{act}. No gradient,
Q value, hidden state, or policy parameter is exposed to the explainer. This
supports the audit of rule-based, value-based, policy-gradient, or hierarchical
responders through the same interface.

\section{Experimental Design}

\subsection{Process and Response Policy}

The independent evaluation environment contains three coupled tanks, four
flows, three PLC assets, a historian, an engineering workstation, and separated
enterprise and control access points. It follows the type of closed-loop
evaluation enabled by water-treatment testbeds such as SWaT~\cite{mathur2016swat}
without using its recorded data. Measurement replay changes the delivered level
while retaining the hidden process state; actuator override changes inlet and
outlet behavior.

The audited policy is a flat tabular Q-learning responder with exponentially
weighted evidence memory. It selects monitor, inspect, restrict, or backup.
Response intensity reduces the active attack effect while imposing an
availability cost. The policy is frozen before counterfactual generation. This
baseline supplies a sequential learned decision to explain; it is not presented
as a new response algorithm.

\subsection{Baselines and Metrics}

Four generators are compared: random valid edits, direct feature-space
counterfactuals, cyber-valid search without physical-impact ranking, and the
complete cyber-physical method. Target validity is the fraction that selects a
weaker response at the audit instant. Cyber feasibility is acceptance by the
independent campaign validator. Physical validity requires a complete bounded
trajectory. Further outcomes are edit cost, impact gain, simulator evaluations,
wall-clock time, and explanation stability across nearby initial conditions.

Selected counterfactuals are then replayed under a smaller integration step,
measurement quantisation, additional integration substeps, and perturbed process
coefficients. Counterfactual transfer validity records whether both the weaker
decision and positive impact gap remain. Training, development, and evaluation
seeds are disjoint. Confirmatory results will be reported over independent seeds
using medians, interquartile ranges, and paired bootstrap intervals.

\InputIfFileExists{generated/confirmatory_table.tex}{}{}

\section{Engineering Pilot and Planned Evaluation}

The current software artifact implements the complete path from a validated
campaign through policy training, frozen-policy execution, counterfactual beam
search, and generated result files. The generated engineering table, when
present, demonstrates executable integration across independent policy seeds.
It is excluded from confirmatory claims. The final evaluation will replace it
with the frozen baseline comparison, feasibility analysis, ablations, and
higher-fidelity transfer.

The primary expected comparison concerns validity rather than raw decision-flip
rate. An unconstrained method may flip decisions frequently while returning
observations that cannot be generated by any valid campaign. The proposed
method is expected to trade some sparsity and search time for reproducible
cyber-physical validity. This trade-off will be reported directly, including
all invalid candidates rather than silently removing them.

\section{Limitations and Conclusion}

The method explains one frozen responder and searches a declared campaign
grammar; it cannot expose tactics absent from that grammar. Feasibility is also
conditional on the fidelity and uncertainty range of the process model. The
claim is consequently limited to identifying executable counterfactuals under
the declared cyber and physical model.

This paper formulates explanation of autonomous ICPS response as constrained
campaign red-teaming. By changing evidence through executable events and
replaying every candidate in closed loop, it connects a changed response to a
reproducible physical consequence. The resulting counterfactuals support policy
audit and provide validated hard campaigns for subsequent work on adaptive
attacker--defender interaction.

\bibliographystyle{IEEEtran}
\bibliography{references}
\end{document}
